% ============================================================================
% CHAPTER 12: PRODUCTION NLP AND DEPLOYMENT
% ============================================================================
% Canonical home for tokenization algorithms (BPE, WordPiece, Unigram,
% SentencePiece).  Around it: the LLM application layer -- serving primitives
% and their application consequences, the compression decision, cost
% engineering, multilingual production challenges, and operating an LLM
% feature.  Serving/compression mechanics are canonical in [DL 19] / [DL 14];
% the MLOps lifecycle is [DL 22].
% ============================================================================

\chapter{Production NLP and Deployment}
\label{chap:production_nlp}

\begin{tldr}
This chapter is the program's canonical home for \textbf{tokenization
algorithms}---BPE, WordPiece, Unigram, SentencePiece, byte-level BPE, and the vocabulary
trade-offs that follow---and, around that core, the \textbf{LLM application layer}: the
serving primitives an application team must specify and design prompts around, the
precision-and-compression decision it actually makes, cost engineering (routing, semantic
caching, prefix-cache-friendly prompts, output-length control), the multilingual tokenizer
tax, and what it takes to operate an LLM feature whose quality has no label.  The serving
and compression \emph{mechanics} are canonical in [DL~19] and [DL~14]; the MLOps lifecycle
is [DL~22].  Every concept is grounded in concrete numbers---memory budgets, latency
targets, cost per request---because interviewers at L5+ expect you to reason quantitatively
about deployment tradeoffs, not just recite architectures.
\end{tldr}

\bigskip

% ============================================================================
% SECTION 1: TOKENIZATION ALGORITHMS
% ============================================================================
\section{Tokenization Algorithms}
\label{sec:tokenization_algorithms}

Tokenization---splitting raw text into the discrete units a model consumes---is deceptively
one of the most consequential design decisions in any NLP system.  The choice of tokenizer
determines the vocabulary size, the sequence length (and therefore compute cost), the model's
ability to handle rare words and multiple languages, and even whether the model can perform
character-level tasks like spelling.

\emph{This section is the program's canonical home for tokenization algorithms}: the
derivations, the vocabulary trade-offs, and the comparison table below are the reference the
other volumes point to.  Two companions live elsewhere: the from-scratch BPE
\emph{implementation drill} (train the merges, encode, handle the degenerate cases) is
[DL~28], and the \emph{interaction effects}---how tokenizer choice couples to the
embedding table, the architecture, and downstream quality---are [DL~9], with the
embedding-side diagnostics in [DL~6].  The multilingual consequences of these algorithms
are developed later in this chapter (Section~\ref{sec:tokenizer_tax}).

\subsection{Why Tokenization Matters}

\begin{enumerate}[label=\textbf{\arabic*.}]
    \item \textbf{Vocabulary size.}  Word-level tokenization produces vocabularies of
          hundreds of thousands of entries, most of which are rare.  Subword tokenization
          compresses this to 32K--100K tokens while retaining coverage of arbitrary text.
    \item \textbf{Sequence length.}  Every token costs compute ($O(n^2)$ in attention).
          A tokenizer that produces 30 tokens for the same sentence another produces 10
          tokens for is $9\times$ more expensive in attention FLOPs.
    \item \textbf{OOV handling.}  Word-level tokenizers assign a single \texttt{[UNK]}
          token to unseen words, discarding all information.  Subword tokenizers decompose
          unseen words into known pieces, preserving morphological signal.
    \item \textbf{Multilinguality.}  Tokenizer design directly determines how well a model
          handles non-English languages.  A tokenizer trained primarily on English will
          shatter CJK or Thai text into byte-level fragments, inflating sequence length.
\end{enumerate}

\begin{warningbox}
Interviewers use tokenization questions as a litmus test for production awareness.  A
candidate who cannot explain how BPE works, or who does not know that tokenizer choice
affects multilingual performance, reveals a gap between paper-reading knowledge and
system-building experience.
\end{warningbox}

\subsection{Byte Pair Encoding (BPE)}
\label{sec:bpe}

BPE (Sennrich et al., ACL 2016) is the most widely used subword tokenization algorithm
and the foundation for GPT-2, GPT-3, GPT-4, LLaMA, and Mistral.

\paragraph{Training algorithm.}
\begin{enumerate}
    \item \textbf{Initialize} the vocabulary with all individual characters (or bytes)
          present in the training corpus.
    \item \textbf{Count} the frequency of every adjacent symbol pair in the corpus.
    \item \textbf{Merge} the most frequent pair into a new symbol and add it to the
          vocabulary.
    \item \textbf{Repeat} steps 2--3 for a predefined number of merges (determining
          the final vocabulary size).
\end{enumerate}

\paragraph{Inference (encoding).}
Given a new word, apply the learned merge rules \emph{greedily} in the order they were
learned during training.  The word is first split into characters, then iteratively
merged according to the priority of each rule.

\paragraph{Handling OOV words.}
BPE never produces an \texttt{[UNK]} token (assuming byte-level BPE).  Any unseen word
is decomposed by splitting it into characters (or bytes) and applying the learned merges
in priority order; whatever subwords the merges produce are the output.  (Greedy
longest-match-first lookup is \emph{WordPiece} inference, not BPE.)  For example, a rare
word like ``unparalleledly'' might tokenize as \texttt{["un", "parallel", "ed", "ly"]}.

\begin{keyinsight}[BPE Is a Compression Algorithm]
BPE was originally a data compression algorithm (Gage, 1994).  Its application to NLP is
natural: the most frequent substrings get their own tokens, just as a compression
dictionary assigns short codes to frequent patterns.  This framing helps explain why
vocabulary size is a compression--compute tradeoff: larger vocabulary $=$ shorter
sequences (better compression) but larger embedding tables and softmax layers.
\end{keyinsight}

\subsection{WordPiece}
\label{sec:wordpiece}

WordPiece (Schuster and Nakajima, 2012; used in BERT, DistilBERT) is structurally
similar to BPE but differs in its merge criterion.

\paragraph{Key difference from BPE.}
BPE merges the pair with the highest \emph{frequency}.  WordPiece merges the pair that
maximizes the \emph{likelihood} of the training data under a unigram language model:
\[
    \text{score}(x, y) = \frac{\text{freq}(xy)}{\text{freq}(x) \cdot \text{freq}(y)}
\]
This prefers merges that create tokens which are more than the sum of their parts---pairs
that co-occur more than expected by chance.  In practice, the outputs of BPE and WordPiece
are similar, but WordPiece tends to produce slightly more linguistically meaningful tokens.

\paragraph{Special prefix.}
WordPiece marks continuation tokens with a \texttt{\#\#} prefix.  The word
``playing'' might tokenize as \texttt{["play", "\#\#ing"]}, explicitly indicating that
\texttt{\#\#ing} is not a word start.

\subsection{Unigram Language Model}
\label{sec:unigram}

The Unigram model (Kudo, EMNLP 2018) takes the opposite approach to BPE:

\paragraph{Training algorithm.}
\begin{enumerate}
    \item \textbf{Initialize} with a large seed vocabulary (e.g., all substrings up to
          a maximum length, or the top 1M most frequent substrings).
    \item \textbf{Estimate} token probabilities using EM: the E-step finds the most
          probable segmentation of each word; the M-step updates token probabilities
          from these segmentations.
    \item \textbf{Prune} tokens whose removal causes the least increase in corpus
          likelihood.  Remove a fixed percentage (e.g., 20\%) of the vocabulary.
    \item \textbf{Repeat} steps 2--3 until the target vocabulary size is reached.
\end{enumerate}

\paragraph{Inference (encoding).}
Unigram finds the \emph{most probable} segmentation using the Viterbi algorithm.  This
is a key advantage over BPE: Unigram has a principled probabilistic model over
segmentations, while BPE's greedy merge is a heuristic.

\paragraph{Used by:} T5, XLNet, ALBERT (all via SentencePiece).  Note that LLaMA also
uses SentencePiece, but with the \emph{BPE} algorithm, not Unigram.

\subsection{SentencePiece and Byte-Level BPE}

\paragraph{SentencePiece} (Kudo and Richardson, EMNLP 2018) is a tokenizer library
(not an algorithm) that implements both BPE and Unigram.  Its key feature is
\emph{language agnosticism}: it treats the input as a raw stream of Unicode text with no
pre-tokenization (no whitespace splitting, no language-specific rules), with optional
byte-fallback for characters outside the vocabulary.  This makes it
essential for multilingual models.  Whitespace is treated as a regular character
(represented as \texttt{\textunderscore}).

\paragraph{Byte-level BPE} (used in GPT-2+) operates on raw UTF-8 bytes rather than
Unicode characters.  Since there are only 256 possible byte values, the base vocabulary
is tiny and \emph{every possible text} can be represented---no \texttt{[UNK]} tokens,
ever.  This is critical for handling arbitrary scripts, code, and mixed-language text.

\subsection{Tokenization for Code}

Code has properties that natural language tokenizers handle poorly: significant whitespace
(Python indentation), long identifiers (\texttt{getUserAccountBalance}), operators, and
structured syntax.  Production approaches include:
\begin{itemize}
    \item \textbf{Whitespace normalization:} Replace runs of spaces with special tokens
          (e.g., Codex uses tokens for 1--24 spaces).
    \item \textbf{CamelCase/snake\_case splitting:} Some tokenizers split identifiers
          at case boundaries or underscores before applying BPE.
    \item \textbf{Larger vocabulary:} Code-focused models often use larger vocabularies
          (e.g., 100K+) to represent common code patterns as single tokens.
\end{itemize}

\subsection{Vocabulary Size Tradeoffs}

\begin{table}[H]
\centering
\small
\begin{tabularx}{\textwidth}{@{}l c c X@{}}
\toprule
\textbf{Vocab Size} & \textbf{Seq.\ Length} & \textbf{Embedding Params} & \textbf{Models} \\
\midrule
$\sim$30--32K & Longer   & 130M (at $d{=}4096$) & BERT (30,522), LLaMA 1/2 (32K) \\
$\sim$50K  & Moderate & 200M (at $d{=}4096$) & GPT-2/GPT-3 (50,257) \\
$\sim$100--128K & Shorter  & 400--520M (at $d{=}4096$) & GPT-4 ($\sim$100K), LLaMA 3 (128K) \\
$\sim$256K & Shortest & 1B (at $d{=}4096$) & Gemini, Gemma \\
\bottomrule
\end{tabularx}
\caption{Vocabulary size tradeoffs.  Larger vocabularies reduce sequence length (saving
compute in attention) but increase the embedding table size.  Modern models trend toward
larger vocabularies as the attention savings outweigh the embedding cost.}
\label{tab:vocab_size}
\end{table}

\subsection{Comparison of Tokenization Methods}

\begin{table}[H]
\centering
\small
\begin{tabularx}{\textwidth}{@{}l l l l X@{}}
\toprule
\textbf{Method} & \textbf{Merge Criterion} & \textbf{Direction} & \textbf{UNK?} & \textbf{Used By} \\
\midrule
BPE              & Most frequent pair        & Bottom-up & No (byte-level)  & GPT-2/3/4, LLaMA, Mistral \\
WordPiece        & Max likelihood gain       & Bottom-up & Possible         & BERT, DistilBERT \\
Unigram          & Min likelihood loss        & Top-down  & No (with bytes)  & T5, XLNet, ALBERT \\
SentencePiece    & Library (BPE or Unigram)  & Either    & No               & LLaMA, T5, mBART \\
Byte-level BPE   & Frequency (on bytes)       & Bottom-up & Never            & GPT-2+, all modern LLMs \\
\bottomrule
\end{tabularx}
\caption{Comparison of subword tokenization methods.}
\label{tab:tokenizer_comparison}
\end{table}

\subsection{TikZ: BPE Merge Process}

\begin{figure}[H]
\centering
\begin{tikzpicture}[
    node distance=0.6cm,
    tok/.style={rectangle, draw=deepblue, fill=lightblue, minimum width=0.7cm,
                minimum height=0.55cm, font=\footnotesize\ttfamily, rounded corners=2pt,
                thick, inner sep=3pt},
    mergetok/.style={rectangle, draw=deepgreen!80, fill=green!8, minimum width=0.7cm,
                minimum height=0.55cm, font=\footnotesize\ttfamily, rounded corners=2pt,
                thick, inner sep=3pt},
    arr/.style={-{Stealth[length=4pt]}, thick, deepblue!60},
    steplbl/.style={font=\small\sffamily\bfseries, text=deepblue},
    mergelbl/.style={font=\scriptsize\sffamily, text=deepgreen!80!black}
]

% Step 0: Characters
\node[steplbl] (s0) at (0, 0) {Init:};
\node[tok, right=0.3cm of s0] (c1) {l};
\node[tok, right=0.08cm of c1] (c2) {o};
\node[tok, right=0.08cm of c2] (c3) {w};
\node[tok, right=0.08cm of c3] (c4) {e};
\node[tok, right=0.08cm of c4] (c5) {s};
\node[tok, right=0.08cm of c5] (c6) {t};

% Step 1: Merge (e,s) -> es
\node[steplbl, below=0.7cm of s0] (s1) {Merge 1:};
\node[tok, right=0.3cm of s1] (m1a) {l};
\node[tok, right=0.08cm of m1a] (m1b) {o};
\node[tok, right=0.08cm of m1b] (m1c) {w};
\node[mergetok, right=0.08cm of m1c] (m1d) {es};
\node[tok, right=0.08cm of m1d] (m1e) {t};
\node[mergelbl, right=0.3cm of m1e] {(e, s) $\to$ es};

% Step 2: Merge (es,t) -> est
\node[steplbl, below=0.7cm of s1] (s2) {Merge 2:};
\node[tok, right=0.3cm of s2] (m2a) {l};
\node[tok, right=0.08cm of m2a] (m2b) {o};
\node[tok, right=0.08cm of m2b] (m2c) {w};
\node[mergetok, right=0.08cm of m2c] (m2d) {est};
\node[mergelbl, right=0.3cm of m2d] {(es, t) $\to$ est};

% Step 3: Merge (l,o) -> lo
\node[steplbl, below=0.7cm of s2] (s3) {Merge 3:};
\node[mergetok, right=0.3cm of s3] (m3a) {lo};
\node[tok, right=0.08cm of m3a] (m3b) {w};
\node[tok, right=0.08cm of m3b] (m3c) {est};
\node[mergelbl, right=0.3cm of m3c] {(l, o) $\to$ lo};

% Step 4: Merge (lo,w) -> low
\node[steplbl, below=0.7cm of s3] (s4) {Merge 4:};
\node[mergetok, right=0.3cm of s4] (m4a) {low};
\node[tok, right=0.08cm of m4a] (m4b) {est};
\node[mergelbl, right=0.3cm of m4b] {(lo, w) $\to$ low};

% Step 5: Merge (low,est) -> lowest
\node[steplbl, below=0.7cm of s4] (s5) {Merge 5:};
\node[mergetok, right=0.3cm of s5] (m5a) {lowest};
\node[mergelbl, right=0.3cm of m5a] {(low, est) $\to$ lowest};

% Arrows between steps
\draw[arr] ([yshift=-0.15cm]c3.south) -- ([yshift=0.15cm]m1c.north);
\draw[arr] ([yshift=-0.15cm]m1d.south) -- ([yshift=0.15cm]m2d.north);
\draw[arr] ([yshift=-0.15cm]m2a.south) -- ([yshift=0.15cm]m3a.north);
\draw[arr] ([yshift=-0.15cm]m3a.south) -- ([yshift=0.15cm]m4a.north);
\draw[arr] ([yshift=-0.15cm]m4a.south) -- ([yshift=0.15cm]m5a.north);

\end{tikzpicture}
\caption{BPE merge process for the word ``lowest.''  Starting from individual characters,
the algorithm iteratively merges the most frequent adjacent pair.  The merge order
becomes the encoding rule: at inference time, the same merges are applied greedily
to any input text.}
\label{fig:bpe_merge}
\end{figure}


% ============================================================================
% SECTION 2: SERVING AND EFFICIENCY (APPLICATION VIEW)
% ============================================================================
\section{Serving and Efficiency: The Application View}
\label{sec:model_serving}

Serving is latency-optimized where training is throughput-optimized, and the
machinery that makes it work is canonical in \textbf{[DL~19]}: the KV cache
and its memory arithmetic, MLA, PagedAttention, prefix caching and
RadixAttention, continuous batching, the prefill/decode split and the
TTFT--throughput frontier, speculative decoding, structured and constrained
decoding, quantization for inference, serving parallelism and MoE serving,
and capacity planning for reasoning models.  The compression methods
themselves---PTQ and QAT, GPTQ and AWQ, pruning, distillation---are canonical
in \textbf{[DL~14]}, and the training-side techniques a serving discussion
tends to collect by accident---mixed precision, gradient checkpointing---are
canonical in \textbf{[DL~15]}.

This section keeps the half that belongs to the \emph{application} team---the
team that usually does not own the serving stack but must specify it, budget
for it, and design prompts and traffic around it: which primitives to demand,
what each one buys, and how the cost arithmetic falls out.

\subsection{The Serving Primitives and What Each One Buys}
\label{sec:kv_cache}
\label{sec:continuous_batching}
\label{sec:paged_attention}
\label{sec:prefix_caching}
\label{sec:parallelism}
\label{sec:speculative_decoding}

Generation runs in two phases, and every serving decision traces back to the
split.  \textbf{Prefill} processes all $n$ prompt tokens in parallel and is
compute-bound; it sets time-to-first-token.  \textbf{Decode} emits one token
at a time, re-reading the whole cache per token, and is
memory-bandwidth-bound; it sets inter-token latency.  The cache itself costs
$M_{\mathrm{KV}} = 2 \cdot L \cdot H_{\mathrm{kv}} \cdot S \cdot d_h \cdot b$
bytes per request---linear in context length, which is why long prompts are a
memory decision and not only a token-price decision (worked numbers, GQA/MLA
reductions, and the full derivation: [DL~19]).
% label preserved from the pre-dedup KV-cache figure (Phase D); the diagram is canonical in [DL 19]
\label{fig:kv_cache}

\begin{table}[H]
\centering
\small
\begin{tabularx}{\textwidth}{@{}l X X@{}}
\toprule
\textbf{Primitive} & \textbf{Mechanism in one line} & \textbf{What it buys, and what it means for your application} \\
\midrule
KV cache & Store per-position keys/values so each new token attends to cached state instead of recomputing it & Turns $O(t)$ work per token into $O(1)$; makes decode bandwidth-bound. Your context length is a memory budget \\
GQA / MLA & Fewer KV heads, or a compressed latent cache & $4$--$8\times$ smaller cache, so more concurrent requests and longer contexts per GPU \\
Continuous batching & Iteration-level scheduling: finished sequences leave the batch, new ones join at any decode step & $2$--$8\times$ throughput over static batching. Your p99 now depends on co-tenant traffic, so measure under realistic load \\
PagedAttention & Block-allocated cache with a page table, no contiguous reservation & Memory utilization $\sim$20\% $\to$ $\sim$95\%, larger batches, no over-allocation for unknown output length \\
Prefix caching & Reuse the computed KV of a shared prefix across requests & Cuts prefill cost and TTFT for a shared system prompt. \emph{Design consequence:} put stable content first and per-request content last---one interpolated timestamp or user ID at the top invalidates the prefix for every request \\
Speculative decoding & Small draft model proposes $k$ tokens, target model verifies them in one pass; rejection sampling keeps the output distribution exact & $2$--$3\times$ faster decode with no quality change. Helps latency at moderate load; at saturation it spends compute you needed for throughput \\
Tensor / pipeline parallelism & Split each layer across GPUs (TP) or assign layer ranges to stages (PP) & Fits models that exceed one GPU. TP inside an NVLink node, PP across nodes; both add failure modes to your capacity plan \\
\bottomrule
\end{tabularx}
\caption{Serving primitives an LLM application depends on, and the consequence
each one has for application design.  Mechanisms, memory math, and the
throughput/latency frontier are canonical in [DL~19].}
\label{tab:parallelism}
\end{table}

\begin{keyinsight}[Prompt Design Is Latency Design]
The application team controls one input to the serving stack that dominates
TTFT: prompt length and prompt \emph{stability}.  Prefill cost is linear in
prompt tokens, so a 4,000-token system prompt that could have been 800 tokens
is a $5\times$ prefill bill on every request; and because prefix caching keys
on an exact shared prefix, ordering the prompt as
\texttt{[stable instructions][retrieved context][user turn]} is worth more
than most infrastructure tuning available to you.  This is the cheapest
latency work in the stack and it belongs to the people writing the prompts,
not the people running the GPUs.
\end{keyinsight}

\subsection{Precision and Compression: The Application-Side Decision}
\label{sec:quantization}
\label{sec:distillation}
\label{sec:efficiency}

Three levers shrink a model, and an application team picks among them far more
often than it implements them.  \textbf{Quantization} lowers numerical
precision (FP16 $\to$ FP8/INT8 $\to$ INT4 via GPTQ or AWQ): no retraining,
$2$--$4\times$ memory and a similar speedup, and it is the default first move.
\textbf{Distillation} trains a smaller student on a larger teacher's outputs
or logits: it needs training data and a training loop, but it is the only
lever that changes the architecture, and for a narrow application task a
distilled small model routinely beats a quantized large one at equal cost.
\textbf{Pruning} removes weights or whole heads and layers: structured
pruning yields a plain dense model, while unstructured sparsity needs
hardware support to turn into speed.  They compose---distill, then quantize
the student---and the mechanics of all three are canonical in [DL~14], with
inference-time quantization in [DL~19].

\begin{table}[H]
\centering
\small
\begin{tabularx}{\textwidth}{@{}l X X X@{}}
\toprule
& \textbf{Quantization} & \textbf{Distillation} & \textbf{Pruning} \\
\midrule
\textbf{What changes} & Numerical precision & Architecture (smaller student) & Sparsity / removed heads \\
\textbf{Cost to apply} & Hours (PTQ), no training data & A training run plus teacher outputs & Minutes to a short fine-tune \\
\textbf{Typical gain} & $2$--$4\times$ memory, $2$--$3\times$ speed & $2$--$6\times$ speed at 90--97\% of teacher & $1.5$--$3\times$, hardware permitting \\
\textbf{Best when} & You need the same model, cheaper & The task is narrow and you own data & You need a dense model slightly smaller \\
\bottomrule
\end{tabularx}
\caption{The three compression levers as an application team experiences them.
Method-level detail (GPTQ vs.\ AWQ, QAT, structured vs.\ unstructured pruning,
distillation objectives) is canonical in [DL~14].}
\label{tab:quantization_tradeoffs}
\label{tab:compression_comparison}
\end{table}

\begin{warningbox}
Quantization quality loss varies dramatically across models, tasks, and
evaluation benchmarks.  A 4-bit quantized model might show only 1\%
degradation on MMLU but 10\% on a math reasoning benchmark, and long-context
or tool-calling behavior often degrades before any headline benchmark moves.
\emph{Always evaluate quantized models on your specific task, with your own
prompts, before deploying.}  The ``1--3\% quality loss'' numbers widely cited
are averages over easy benchmarks and are not a warranty for your use case.
\end{warningbox}

\subsection{Cost Engineering for LLM Applications}
\label{sec:cost_engineering}

Application cost is an arithmetic identity before it is an optimization
problem: cost per request $=$ (prompt tokens $\times$ input price) $+$ (output
tokens $\times$ output price), summed over every model call the request makes,
including the ones users never see---guardrail classifiers, rerankers,
LLM-judge sampling, and retries.  Two consequences follow immediately: a
feature's cost is set by its \emph{prompt architecture}, not by its model
choice alone, and the tokenizer tax (Section~\ref{sec:tokenizer_tax}) silently
multiplies both terms for non-English users.

\begin{itemize}
    \item \textbf{Model routing.}  Classify the query and send easy traffic to
          a small model, hard traffic to a large one.  A well-tuned router
          cuts average cost $5$--$10\times$ while holding quality on the hard
          slice; the engineering risk is the router itself, which needs its
          own eval set and a fail-open path to the large model.
    \item \textbf{Semantic caching.}  Return a stored response when a new
          query's embedding is within a threshold of a cached one.  Excellent
          for FAQ-shaped traffic; dangerous when near-duplicate phrasings
          carry different intent (``how do I cancel'' vs.\ ``how do I cancel
          \emph{without a fee}''), when responses are personalized, or when
          the underlying data changes.  Treat the threshold as a
          precision/recall decision, scope keys by user and tenant, and set a
          TTL tied to how fast the source data moves.
    \item \textbf{Prefix and prompt caching.}  The billing-visible form of the
          serving primitive above: providers charge less for cached prefix
          tokens, so prompt \emph{stability} converts directly into money.
    \item \textbf{Prompt compression.}  Shorten system prompts, summarize
          conversation history beyond a horizon, and retrieve fewer, better
          chunks rather than more ([NLP~9]).  Prefill is linear in tokens, so
          this is the one lever that cuts latency and cost together.
    \item \textbf{Output-length control.}  Output tokens are priced higher and
          are serialized through decode; a response-format instruction and a
          \texttt{max\_tokens} ceiling are cost controls, not cosmetics.
    \item \textbf{Batch and off-peak tiers.}  Anything that does not need an
          interactive answer---backfills, evaluation runs, bulk
          classification---belongs on a batch endpoint at a fraction of the
          interactive price.
\end{itemize}


% ============================================================================
% SECTION 3: MULTILINGUAL PRODUCTION CHALLENGES
% ============================================================================
\section{Multilingual Production Challenges}
\label{sec:multilingual_production}

Deploying NLP systems globally introduces challenges that monolingual (English-focused)
development obscures.

\subsection{The Tokenizer Tax}
\label{sec:tokenizer_tax}

Tokenizers trained primarily on English text produce dramatically more tokens for
non-English and non-Latin-script languages.  This ``tokenizer tax'' has cascading
effects:

\begin{itemize}
    \item \textbf{Effective context window is shorter.}  If Thai text tokenizes at
          $3\times$ the rate of English, a 128K-token context window holds
          $\sim$43K tokens worth of Thai content.
    \item \textbf{Inference is more expensive.}  More tokens $=$ more compute
          ($O(n^2)$ attention) and more KV-cache memory.
    \item \textbf{API costs increase.}  Token-based pricing penalizes non-English users.
    \item \textbf{Generation is slower.}  More tokens to generate for the same semantic
          content.
\end{itemize}

\paragraph{Concrete example.}
The sentence ``Artificial intelligence is transforming the world'' is approximately:
\begin{itemize}
    \item \textbf{English:} 7 tokens (GPT-4 tokenizer)
    \item \textbf{Chinese:} 11 tokens ($1.6\times$)
    \item \textbf{Japanese:} 14 tokens ($2.0\times$)
    \item \textbf{Thai:} 25 tokens ($3.6\times$)
    \item \textbf{Hindi:} 30 tokens ($4.3\times$)
\end{itemize}

\paragraph{Mitigation strategies.}
\begin{itemize}
    \item Train tokenizers on balanced multilingual corpora (equal representation of
          languages).
    \item Increase vocabulary size (100K+ tokens) to include more non-English subwords.
    \item Use byte-level BPE to ensure universal coverage.
    \item Evaluate fertility (tokens per word) across target languages as a tokenizer
          quality metric.
\end{itemize}

\subsection{Cross-Lingual Transfer}

Multilingual models (mBERT, XLM-R, NLLB-200) enable training on high-resource languages
and deploying on low-resource languages.  This works because shared subword vocabularies
and shared model parameters force the model to learn \emph{language-agnostic}
representations.

\paragraph{When cross-lingual transfer fails.}
\begin{itemize}
    \item Script mismatch: Latin-script training data transfers poorly to Devanagari or
          Thai script without explicit cross-script signal.
    \item Morphological mismatch: English (analytic) to Turkish (agglutinative)
          transfer loses morphological structure.
    \item Domain mismatch: legal or medical terminology is language-specific and does
          not transfer.
\end{itemize}

\subsection{Code-Switching}

Real-world text frequently mixes languages: ``Let's meet at the caf\'{e}, va bene?''
Code-switched text breaks assumptions of both language detection and monolingual models.
Production approaches include multilingual models (which handle mixed-language input
natively), language-aware tokenization, and training on code-switched data.


% ============================================================================
% SECTION 4: OPERATING AN LLM APPLICATION
% ============================================================================
\section{Operating an LLM Application}
\label{sec:mlops_nlp}

The MLOps lifecycle---the drift taxonomy and its detection tests, monitoring
architecture, retraining triggers and cadence, CI/CD for models, deployment
patterns (shadow, canary, blue-green), feature stores and training--serving
skew---is canonical in \textbf{[DL~22]}, and the statistics of the A/B tests
that gate those rollouts are canonical in \textbf{[CML~6]}.  None of it is
repeated here.  What follows is what changes when the artifact you operate is
an LLM feature rather than a classifier.

\subsection{What to Watch That a Classifier Never Made You Watch}

\paragraph{Latency has two clocks.}  \textbf{TTFT} (prefill-dominated; target
$<$500\,ms interactive, $<$100\,ms for search-like surfaces) and
\textbf{inter-token latency} (target $<$50\,ms, i.e.\ 20+ tokens/s, for
comfortable reading) are separate SLOs with separate fixes: TTFT is prompt
length and prefix-cache hit rate, ITL is decode and batch pressure.  A single
end-to-end p99 hides both.  Alert on each, and watch \textbf{output-length
distribution} as a leading indicator---a prompt change that makes responses
40\% longer is a latency and cost regression that no accuracy metric shows.

\paragraph{Quality has no label.}  There is no ground truth arriving
downstream, so quality monitoring is assembled from three sources: implicit
signals (regeneration clicks, copy rates, abandonment, thumbs), automated
judging over sampled traffic (LLM-as-judge, factuality and toxicity
classifiers---methodology in [NLP~10]), and a \textbf{golden set} that every
candidate model, prompt, and decoding-parameter change must pass before it
ships.  Track \textbf{refusal rate} and \textbf{cache hit rate} alongside
them: both move sharply on changes nobody expected to affect them.

\paragraph{Drift arrives as text.}  The NLP-specific drift signals are input
length distribution (sudden shifts often mean bot traffic or a new client),
input token-frequency divergence from the reference window, embedding
centroid and spread, and per-language traffic mix.  \textbf{Slice every metric
by language}: the tokenizer tax (Section~\ref{sec:tokenizer_tax}) means
non-English users already pay more tokens for the same request, so their
latency and cost regress first and their quality regresses invisibly if the
dashboard is English-weighted.

\subsection{Prompt and Version Management}

The deployed artifact of an LLM feature is not a model file.  It is a tuple:
\emph{model version + tokenizer + prompt template + decoding parameters +
tool/function schemas + retrieval index version}.  Any one of them can change
behavior on its own, so pin them together, version them together, and record
the tuple with every logged request---an incident you cannot attribute to a
specific tuple is an incident you cannot fix.  Prompts in particular are
code: they belong in the repository, in review, with an eval suite that runs
in CI and a canary rollout by prompt version, not in a spreadsheet or an
untracked console.  Keep instant rollback for every element of the tuple, and
remember that a provider-side model update is a version change you did not
make---pin model versions explicitly and re-run the golden set when a pin
expires.

\subsection{Feedback Loops}

\begin{enumerate}
    \item \textbf{Implicit feedback:}  Behavioral signals---regenerate clicks,
          copy-paste rates, session abandonment, follow-up rephrasings.  Cheap
          and plentiful, but confounded by UI and position effects.
    \item \textbf{Explicit feedback:}  Thumbs up/down, ratings, user
          corrections.  Sparse, biased toward extremes, and the only signal
          that names \emph{why}.
    \item \textbf{Active learning:}  Route low-confidence or
          high-disagreement traffic to human annotators and fold the results
          into the golden set and the training data.
    \item \textbf{Preference loops:}  Production preference data feeds reward
          models and DPO-style updates; the alignment machinery it feeds is
          canonical in [NLP~8].  Guard the loop: a model trained on its own
          accepted outputs narrows over time, so hold out a slice of
          unfiltered traffic to measure what the loop is quietly removing.
\end{enumerate}


% ============================================================================
% SECTION 5: INTERVIEW QUESTIONS
% ============================================================================
\section{Interview Questions}
\label{sec:production_interview_questions}

Several questions below drill serving and compression mechanics whose canonical
treatment is [DL~19] and [DL~14]; they are kept here because an NLP application
engineer is expected to produce these numbers on demand, and the answers are written
from the application side.  [DL~19] carries a separate and larger bank of
serving-design questions---work both if serving is your interview surface.

% --- Question 1: Serving at Scale ---
\begin{interviewq}
\textbf{Q: How would you serve an LLM at 10K QPS with less than 500ms time-to-first-token?}

\badgeLsix{L6} \quad \badgetype{System Design} \quad \badgefreq{\freqhigh}

\stronganswer
This requires reasoning across the full serving stack.  10K QPS with 500ms TTFT is
extremely demanding---let me state assumptions first, then work the token budget
honestly.  \textbf{Assumptions:} 7B-class model with GQA ($H_{\text{kv}}{=}8$),
1K-token prompts, $\sim$200-token responses, H100-class GPUs ($\sim$3.35\,TB/s HBM;
$\sim$990 dense BF16 TFLOPS of which $\sim$50\% is realistically sustained, i.e.\
$\sim$500 effective TFLOPS; FP8 roughly doubles that).

\textbf{Model selection and optimization:}  Start with the smallest model that meets
quality requirements.  A 7B model with INT4 weight quantization (AWQ or GPTQ) fits on
a single GPU with ample room for KV-cache.  TTFT is dominated by prefill, which is
\emph{compute}-bound: a 1K-token prompt costs $2N \approx 14$\,GFLOPs per token
$\times$ 1K tokens $= 14$\,TFLOPs, i.e.\ $\sim$30\,ms at 500 effective TFLOPS
($\sim$15\,ms with FP8)---well inside the 500ms budget \emph{per request}.  The hard
part is aggregate throughput, not single-request latency.  If a 70B model is required,
use GQA + quantization + tensor parallelism across a node, and every count below
scales up by roughly $10\times$.

\textbf{Batching and the token budget:}  Continuous batching is essential, but the GPU
count follows from the \emph{token} budget, not the request count.  10K QPS $\times$
1K prompt tokens $=$ 10M prefill tokens/s.  One H100 sustains roughly
$500\,\text{TFLOPS} / 14\,\text{GFLOPs per token} \approx 35\text{K}$ prefill tokens/s
in BF16 ($\sim$70K with FP8)---only $\sim$35--70 QPS of 1K-token prefill per GPU.
Prefill alone therefore needs $\sim$140 (FP8) to $\sim$280 (BF16) GPUs.  Decode adds
10K QPS $\times$ 200 output tokens $=$ 2M tokens/s; at high batch, decode is also
compute-bound at the same per-GPU token rate, adding another $\sim$30--60 GPUs.
Total: on the order of \textbf{200--350 H100s} ($\sim$25--45 8-GPU nodes), plus
headroom for P99 and traffic spikes---an order of magnitude more than a naive
request-count estimate suggests.

\textbf{Infrastructure:}
\begin{itemize}
    \item Use vLLM or TensorRT-LLM as the serving engine (PagedAttention, continuous
          batching, prefix caching).
    \item Prefix caching for shared system prompts eliminates redundant prefill.
    \item Load balance across GPU replicas with least-connections routing.
    \item Horizontal scaling: add replicas linearly to increase throughput.
\end{itemize}

\textbf{Latency optimization:}  Prefix caching reduces TTFT for repeated prefixes.
Speculative decoding reduces inter-token latency (but TTFT is dominated by prefill,
so speculative decoding helps less here).  FP8 on H100 provides $\sim$2$\times$
prefill speedup over FP16.

\textbf{Capacity planning:}  Apply Little's law: each request spends ${\sim}100$\,ms in
prefill and ${\sim}3$\,s in decode (200 tokens at ${\sim}15$\,ms/token at high batch),
so at 10K QPS roughly $10{,}000 \times 3 \approx 30{,}000$ requests are in flight at any
moment.  Per-request KV-cache for the assumed GQA-8 7B model at ${\sim}1.2$K tokens with
INT4 KV is $2 \cdot L \cdot H_{\text{kv}} \cdot S \cdot d_h \cdot b = 2 \times 32 \times
8 \times 1200 \times 128 \times 0.5 \approx 40$\,MB (the same cache with MHA in FP16
would be ${\sim}630$\,MB).  Cluster-wide that is ${\sim}1.2$\,TB of KV---only a few GB
per GPU spread across ${\sim}300$ GPUs.  With PagedAttention this is manageable; in this
regime the binding constraint is prefill compute, not KV memory.

\redflags
\begin{itemize}
    \item Jumping to ``just add more GPUs'' without reasoning about per-GPU capacity.
    \item Not mentioning continuous batching or PagedAttention.
    \item Ignoring the prefill vs.\ decode distinction for TTFT.
    \item Not quantifying anything---giving a purely qualitative answer.
\end{itemize}

\interviewtest{Whether the candidate can design a complete serving system with concrete
numbers, distinguishing prefill-bound TTFT from decode-bound token generation, and
reasoning about GPU-level capacity.}

\goodgreat
\begin{itemize}
    \item \badgeLfive{L5}: Mentions continuous batching, quantization, and prefix
          caching.  Gives rough GPU count estimates.
    \item \badgeLsix{L6}: Quantifies KV-cache memory, prefill latency, and per-GPU
          throughput.  Distinguishes prefill vs.\ decode bottlenecks.  Discusses
          failover and load balancing.
    \item \badgeLseven{L7}: Designs the full system including auto-scaling, model
          routing (small model for easy queries), graceful degradation under load,
          multi-region deployment, and cost analysis per query.
\end{itemize}

\followups{What changes if the average prompt is 32K tokens instead of 1K?  How would
you handle a 10$\times$ traffic spike?  What is the cost per query?}
\end{interviewq}

\vspace{0.5em}

% --- Question 2: KV-Cache ---
\begin{interviewq}
\textbf{Q: Explain the KV-cache.  Why is it necessary for autoregressive generation?
How much memory does it use?}

\badgeLfive{L5} \quad \badgetype{Conceptual} \quad \badgefreq{\freqhigh}

\stronganswer
In autoregressive generation, each new token attends to all previous tokens.  Without
caching, generating token $t$ requires recomputing the key and value projections for
all $t-1$ previous tokens---redundant work that grows quadratically over the full
sequence.

The KV-cache stores the key and value tensors from all previous positions so they are
computed only once.  Generation has two phases: \textbf{prefill} (process the entire
prompt in parallel, building the initial KV-cache---compute-bound) and \textbf{decode}
(generate one token at a time, appending to the cache---memory-bandwidth-bound).

\textbf{Memory calculation:}
$M_{\text{KV}} = 2 \times L \times H_{\text{kv}} \times S \times d_h \times b$.
For LLaMA-2 7B ($L{=}32$, $H_{\text{kv}}{=}32$, $d_h{=}128$) at 2K context in FP16:
$2 \times 32 \times 32 \times 2048 \times 128 \times 2 = 1.07$\,GB per request.  With
GQA ($H_{\text{kv}}{=}8$): 268\,MB.  This is why GQA is critical for production---it
reduces KV-cache memory by $4\times$ with minimal quality loss.

The KV-cache is the main memory bottleneck at inference time.  For long contexts (128K+
tokens), it can exceed the model weights in memory.  This is why PagedAttention (virtual
memory for KV-cache) and KV-cache quantization are active areas of optimization.

\redflags
\begin{itemize}
    \item Not knowing what the KV-cache stores or why it exists.
    \item Confusing prefill and decode phases.
    \item Inability to estimate memory usage with concrete numbers.
    \item Thinking KV-cache is optional or a ``nice to have.''
\end{itemize}

\interviewtest{Whether the candidate understands the fundamental inference mechanics of
autoregressive transformers and can reason quantitatively about memory.}

\goodgreat
\begin{itemize}
    \item \badgeLfive{L5}: Explains the two-phase structure, gives the memory formula,
          and computes a concrete example.
    \item \badgeLsix{L6}: Discusses GQA/MQA impact on KV-cache size, PagedAttention,
          prefix caching, and the distinction between compute-bound prefill and
          memory-bandwidth-bound decode.
    \item \badgeLseven{L7}: Discusses KV-cache quantization (INT8/INT4 for K/V tensors),
          KV-cache compression via attention sinks, and the interaction between
          KV-cache size and continuous batching capacity.
\end{itemize}

\followups{What is PagedAttention?  How does GQA reduce KV-cache size?  What happens to
KV-cache memory at 128K context?  How does continuous batching interact with KV-cache
management?}
\end{interviewq}

\vspace{0.5em}

% --- Question 3: Speculative Decoding ---
\begin{interviewq}
\textbf{Q: What is speculative decoding?  Why does it produce the same distribution as
standard autoregressive decoding?}

\badgeLsix{L6} \quad \badgetype{Conceptual} \quad \badgefreq{\freqmed}

\stronganswer
Speculative decoding accelerates autoregressive generation by using a small, fast
\emph{draft model} to propose candidate tokens that a large \emph{target model} verifies
in parallel.

The algorithm: (1) the draft model generates $k$ candidate tokens autoregressively (cheap:
$\sim$2ms per token for a 1B model).  (2) The target model processes all $k$ candidates
in a single forward pass (parallel, like prefill).  (3) At each position, compare the
draft model's token probability $q(x)$ to the target model's probability $p(x)$.  If
$q(x) \leq p(x)$, accept.  If $q(x) > p(x)$, accept with probability $p(x)/q(x)$,
reject otherwise.  On rejection, sample a correction from the residual distribution
$\max(0, p(x) - q(x))$ (normalized).

\textbf{Why the distribution is preserved:}  The accept/reject scheme is a form of
rejection sampling.  For any token $x$, the probability of it being accepted and output
is exactly $p(x)$:  If $q(x) \leq p(x)$, the token is always accepted (probability
$q(x)$ of being drafted $\times$ probability 1 of acceptance).  If $q(x) > p(x)$, it is
accepted with probability $p(x)/q(x)$ (so $q(x) \cdot p(x)/q(x) = p(x)$).  The
correction sampling handles the rejection case.  The net result is that every output
token is sampled exactly from $p(\cdot)$.

The speedup comes from the asymmetry between generation and verification: generating
$k$ tokens costs $k$ sequential forward passes, but verifying $k$ tokens costs only
one parallel forward pass.

\redflags
\begin{itemize}
    \item Claiming speculative decoding introduces quality loss.
    \item Not explaining the accept/reject mechanism.
    \item Confusing speculative decoding with beam search or other search methods.
    \item Not understanding why verification is cheaper than generation.
\end{itemize}

\interviewtest{Understanding of autoregressive generation mechanics, and whether the
candidate can reason about the mathematical guarantee of distribution preservation.}

\goodgreat
\begin{itemize}
    \item \badgeLfive{L5}: Correctly describes the draft-verify-accept/reject pipeline
          and states that the output distribution is preserved.
    \item \badgeLsix{L6}: Explains the rejection sampling argument formally, quantifies
          acceptance rates and speedup, discusses draft model selection.
    \item \badgeLseven{L7}: Discusses variants (Medusa, EAGLE, self-speculative), the
          theoretical speedup bound as a function of acceptance rate, and practical
          deployment considerations (batch interaction, variable acceptance rates).
\end{itemize}

\followups{How do you choose the draft model?  What is the theoretical maximum speedup?
How does speculative decoding interact with batching?  What is Medusa?}
\end{interviewq}

\vspace{0.5em}

% --- Question 4: Quantization Comparison ---
\begin{interviewq}
\textbf{Q: Compare GPTQ, AWQ, and INT8 quantization.  When would you use each?}

\badgeLsix{L6} \quad \badgetype{Trade-off} \quad \badgefreq{\freqmed}

\stronganswer
All three are post-training quantization methods, but they differ in precision,
mechanism, and quality-speed tradeoff.

\textbf{INT8 (e.g., LLM.int8(), SmoothQuant):}  Quantizes weights to 8-bit integers.
SmoothQuant migrates quantization difficulty from activations to weights by mathematically
equivalent channel-wise scaling.  $\sim$$2\times$ memory reduction, $\sim$$2\times$ speedup,
$<$1\% quality loss on nearly all benchmarks.  \emph{Use when:} you need a safe,
reliable compression with minimal risk---the conservative choice.

\textbf{GPTQ (4-bit):}  Layer-wise quantization using second-order (Hessian) information.
Quantizes each weight column while compensating for the error in remaining columns.
$\sim$$4\times$ memory reduction, $\sim$$3\times$ speedup, 1--3\% quality loss.  Slower
to quantize (hours for large models) but high-quality results.  \emph{Use when:} you need
4-bit compression and have time for careful quantization; good for offline preparation.

\textbf{AWQ (4-bit):}  Scales important weight channels (identified by activation
magnitude) before quantization.  Simpler and faster to apply than GPTQ, often matches or
exceeds GPTQ quality.  Same memory/speed profile.  \emph{Use when:} you need 4-bit
compression with fast quantization; increasingly the preferred 4-bit method.

\textbf{Decision framework:}
\begin{enumerate}
    \item If the model fits in memory at INT8 and quality is critical $\to$ INT8.
    \item If the model does not fit at INT8 or you need maximum throughput $\to$ AWQ
          (default 4-bit choice) or GPTQ.
    \item If you are on H100+ hardware $\to$ consider FP8 as a drop-in replacement for
          INT8 with potentially better quality.
\end{enumerate}

\redflags
\begin{itemize}
    \item Treating all quantization methods as interchangeable.
    \item Not knowing the mechanism behind GPTQ or AWQ (just ``it makes things smaller'').
    \item Not mentioning that quality loss varies by task and must be evaluated.
\end{itemize}

\interviewtest{Whether the candidate understands quantization at a mechanistic level
and can make principled deployment decisions.}

\goodgreat
\begin{itemize}
    \item \badgeLfive{L5}: Correctly describes all three methods and their tradeoffs.
    \item \badgeLsix{L6}: Explains the mechanisms (Hessian compensation for GPTQ,
          activation awareness for AWQ), discusses per-task evaluation necessity, and
          mentions FP8 as an emerging option.
    \item \badgeLseven{L7}: Discusses quantization of KV-cache separately from weights,
          mixed-precision strategies (4-bit weights + 8-bit KV-cache), and the
          interaction between quantization and other optimizations (speculative decoding,
          tensor parallelism).
\end{itemize}

\followups{How does GPTQ use the Hessian?  Why does AWQ focus on activation magnitude?
Can you quantize the KV-cache separately?  What is FP8 and when would you use it?}
\end{interviewq}

\vspace{0.5em}

% --- Question 5: Latency Reduction ---
\begin{interviewq}
\textbf{Q: How would you reduce the latency of an NLP inference system by 5$\times$?}

\badgeLfive{L5} \quad \badgetype{System Design} \quad \badgefreq{\freqhigh}

\stronganswer
A $5\times$ reduction requires stacking multiple optimizations.  I would approach this
systematically, starting with the highest-impact changes.

\textbf{1. Quantization ($2$--$4\times$):}  Move from FP16 to INT4 (AWQ/GPTQ).  This
alone can provide $\sim$$3\times$ speedup for memory-bandwidth-bound decoding, as the
model weights are $4\times$ smaller and read $4\times$ faster.

\textbf{2. Distillation or model selection ($2$--$6\times$):}  If the task allows it,
switch from a 70B model to a fine-tuned 7B model.  The smaller model is $\sim$$10\times$
faster, and with task-specific fine-tuning, quality loss can be minimal for narrow tasks.

\textbf{3. Speculative decoding ($2$--$3\times$):}  Reduces inter-token latency (but not
TTFT) by generating multiple tokens per target-model forward pass.

\textbf{4. Model compilation ($1.3$--$2\times$):}  TensorRT-LLM or \texttt{torch.compile}
fuse operations and optimize kernel selection.

\textbf{5. Serving infrastructure:}  Flash Attention (reduces prefill time), prefix
caching (eliminates redundant prefill), continuous batching (improves throughput and
reduces queuing delay).

\textbf{6. Architecture changes:}  Sliding window attention for long contexts, GQA to
reduce KV-cache bandwidth.

In practice, combining quantization ($3\times$) + compilation ($1.5\times$) already
yields $\sim$$4.5\times$.  Adding speculative decoding for decode-heavy workloads pushes
past $5\times$.  If the quality budget allows a smaller model, the gain is even larger.

\redflags
\begin{itemize}
    \item Giving only one technique (e.g., ``just quantize it'').
    \item Not stacking optimizations with multiplicative reasoning.
    \item Not distinguishing between TTFT and ITL optimization.
    \item Suggesting optimizations without estimating their individual contribution.
\end{itemize}

\interviewtest{Systematic optimization thinking---can the candidate identify multiple
independent axes of improvement and reason about their combined effect?}

\goodgreat
\begin{itemize}
    \item \badgeLfive{L5}: Lists 3+ techniques with approximate speedups and explains
          why they stack multiplicatively.
    \item \badgeLsix{L6}: Distinguishes TTFT vs.\ ITL optimizations, discusses
          quality-latency Pareto frontier, includes model selection as an option.
    \item \badgeLseven{L7}: Designs a full optimization roadmap with measurement at each
          stage, discusses diminishing returns, considers cost vs.\ latency tradeoff,
          and mentions model routing as a system-level optimization.
\end{itemize}

\followups{Which of these optimizations affect TTFT vs.\ decode latency?  How would you
measure the quality impact of each?  What if you cannot change the model?}
\end{interviewq}

\vspace{0.5em}

% --- Question 6: BPE Tokenization ---
\begin{interviewq}
\textbf{Q: Explain BPE tokenization.  How is it trained?  How does it handle OOV words?}

\badgeLfive{L5} \quad \badgetype{Conceptual} \quad \badgefreq{\freqhigh}

\stronganswer
Byte Pair Encoding is a subword tokenization algorithm that learns a vocabulary of
variable-length tokens from a training corpus.

\textbf{Training:}  (1) Initialize the vocabulary with all individual characters (or
bytes, in byte-level BPE).  (2) Count the frequency of every adjacent token pair in the
corpus.  (3) Merge the most frequent pair into a new token and add it to the vocabulary.
(4) Repeat steps 2--3 for a predefined number of merges (e.g., 32,000 merges for a 32K
vocabulary).  The ordered list of merge rules \emph{is} the tokenizer.

\textbf{Encoding (inference):}  Given a new string, start with individual characters and
apply the learned merge rules greedily in priority order.  The result is a sequence of
subword tokens.

\textbf{OOV handling:}  BPE \emph{never} produces an unknown token.  Any unseen word is
decomposed into known subword units.  In byte-level BPE (used by GPT-2+), the base
vocabulary is the 256 possible byte values, so \emph{any} byte sequence has a valid
tokenization.  A rare word like ``unparalleledly'' might become
\texttt{["un", "parallel", "ed", "ly"]}---the model retains morphological information
even for words it has never seen as a complete unit.

\textbf{Vocabulary size tradeoff:}  Larger vocabulary $\to$ shorter sequences (less
compute) but larger embedding tables.  32K--100K is typical; modern models trend toward
100K+ because the attention savings from shorter sequences outweigh the embedding cost.

\redflags
\begin{itemize}
    \item Not knowing the training algorithm (just ``it splits words into pieces'').
    \item Saying BPE produces \texttt{[UNK]} tokens.
    \item Confusing BPE with WordPiece or not knowing the difference.
    \item Not mentioning the vocabulary size tradeoff.
\end{itemize}

\interviewtest{Whether the candidate understands tokenization as an algorithm, not just
a preprocessing step---and whether they appreciate its impact on model behavior.}

\goodgreat
\begin{itemize}
    \item \badgeLfive{L5}: Describes the full training and encoding algorithm, explains
          OOV handling, and mentions the vocabulary size tradeoff.
    \item \badgeLsix{L6}: Compares BPE with WordPiece and Unigram, discusses byte-level
          BPE for multilingual handling, and explains why LLMs struggle with
          character-level tasks (tokenization breaks character boundaries).
    \item \badgeLseven{L7}: Discusses the tokenizer tax for non-English languages,
          tokenizer training as a separate optimization problem from model training,
          and the impact of tokenizer design on cross-lingual transfer.
\end{itemize}

\followups{What is the difference between BPE and WordPiece?  Why do LLMs struggle
with counting letters in words?  How does tokenizer choice affect multilingual
performance?}
\end{interviewq}

\vspace{0.5em}

% --- Question 7: Distillation vs Quantization vs Pruning ---
\begin{interviewq}
\textbf{Q: When would you choose knowledge distillation over quantization or pruning to
compress a model?}

\badgeLfive{L5} \quad \badgetype{Trade-off} \quad \badgefreq{\freqmed}

\stronganswer
These three compression techniques are complementary, not competing.  Each has a
different profile:

\textbf{Quantization:}  Reduces numerical precision.  $2$--$4\times$ smaller, minimal
quality loss, no training required (PTQ).  Choose quantization as the \emph{default first
step}---it is the lowest-effort, highest-reward compression.

\textbf{Pruning:}  Removes weights or structures.  Structured pruning (removing heads or
layers) gives real speedup without hardware constraints; unstructured pruning requires
sparse hardware support.  Choose pruning when you need a smaller \emph{dense} model
(structured) or when targeting hardware with sparse acceleration (unstructured).

\textbf{Distillation:}  Trains a new, smaller architecture.  Highest potential
compression ($5$--$10\times$ with architecture change) but requires a full training run.
Choose distillation when: (1) you need $>$$4\times$ compression (beyond what quantization
provides), (2) you want to change the architecture (e.g., 12-layer $\to$ 6-layer), or
(3) you have a strong teacher model and sufficient compute for student training.

\textbf{The practical order:}
\begin{enumerate}
    \item Start with \textbf{quantization} (INT8 or INT4)---immediate gain, zero training.
    \item If more compression is needed, \textbf{prune} unnecessary heads/layers.
    \item If quality is still insufficient at the target size, \textbf{distill} into a
          properly sized student.
    \item Apply quantization to the distilled/pruned model for additional compression.
\end{enumerate}

\redflags
\begin{itemize}
    \item Treating the three as mutually exclusive.
    \item Recommending distillation when quantization would suffice (wasteful).
    \item Not mentioning that distillation requires training compute.
\end{itemize}

\interviewtest{Practical engineering judgment about when to use which compression
technique, and awareness that they compose.}

\goodgreat
\begin{itemize}
    \item \badgeLfive{L5}: Correctly describes all three techniques, gives scenarios
          for each, and notes they can be combined.
    \item \badgeLsix{L6}: Discusses the practical order of application, mentions
          specific methods (GPTQ, DistilBERT), and quantifies typical compression
          ratios for each technique.
    \item \badgeLseven{L7}: Discusses the quality-efficiency Pareto frontier, mentions
          that distillation can use synthetic data from the teacher (no labeled data
          needed), and considers the full lifecycle cost (training cost of distillation
          vs.\ ongoing inference cost savings).
\end{itemize}

\followups{Can you quantize a distilled model?  How does structured pruning differ from
distillation?  What is the cost-benefit analysis of distillation (training cost vs.\
inference savings)?}
\end{interviewq}

\vspace{0.5em}

% --- Question 8: Tokenizer Tax ---
\begin{interviewq}
\textbf{Q: How do you handle the ``tokenizer tax'' when deploying NLP systems for
multilingual users?}

\badgeLsix{L6} \quad \badgetype{System Design} \quad \badgefreq{\freqmed}

\stronganswer
The tokenizer tax refers to the phenomenon where non-English text (especially non-Latin
scripts) is tokenized into significantly more tokens than equivalent English text.  For
example, the same sentence might be 7 tokens in English but 30 tokens in Hindi with a
tokenizer trained primarily on English.  This has four cascading effects: shorter
effective context windows, higher compute cost, higher API cost, and slower generation.

\textbf{Mitigation at the tokenizer level:}
\begin{itemize}
    \item Train the tokenizer on a \emph{balanced} multilingual corpus, ensuring
          non-English languages have adequate representation.  LLaMA 3's tokenizer was
          explicitly redesigned for better multilingual fertility.
    \item Increase vocabulary size (100K--150K tokens) to include more non-English
          subwords.  This reduces fertility for non-English languages at the cost of
          a larger embedding table.
    \item Use byte-level BPE (ensures universal coverage) combined with multilingual
          training data.
    \item Measure \textbf{fertility} (tokens per word) across all target languages as a
          tokenizer evaluation metric.  Flag any language with fertility $>$$2\times$
          the English baseline.
\end{itemize}

\textbf{Mitigation at the system level:}
\begin{itemize}
    \item For API-based systems, consider per-character or per-word pricing rather than
          per-token to avoid penalizing non-English users.
    \item Allocate larger context budgets for languages with higher fertility.
    \item Use language-specific adapter models or language-specific fine-tuning to
          compensate for reduced per-token quality in underrepresented languages.
    \item Monitor quality metrics \emph{per language}, not just on average.
\end{itemize}

\redflags
\begin{itemize}
    \item Not knowing what the tokenizer tax is.
    \item Suggesting ``just use a bigger context window'' without addressing the root cause.
    \item Ignoring the cost and latency implications.
    \item Not mentioning tokenizer training as the primary lever.
\end{itemize}

\interviewtest{Awareness of multilingual deployment challenges and ability to address
them at both the model and system level.}

\goodgreat
\begin{itemize}
    \item \badgeLfive{L5}: Explains the tokenizer tax with concrete examples and
          suggests training on balanced multilingual data.
    \item \badgeLsix{L6}: Discusses vocabulary size tradeoffs, fertility as a metric,
          and system-level mitigations (pricing, context budgets, per-language monitoring).
    \item \badgeLseven{L7}: Discusses cross-lingual transfer failure modes (script
          mismatch, morphological mismatch), the tension between shared vocabulary
          and per-language quality, and strategies like language-adaptive tokenization
          or language mixture training schedules.
\end{itemize}

\followups{How would you evaluate tokenizer quality across 50 languages?  What is the
relationship between tokenizer fertility and downstream quality?  How does tokenizer
design affect cross-lingual transfer?}
\end{interviewq}

\vspace{0.5em}

% --- Question 9: Production Monitoring ---
\begin{interviewq}
\textbf{Q: You have deployed an LLM-powered feature and users report quality has degraded
over the past week, but your automated metrics show no change.  How do you diagnose this?}

\badgeLsix{L6} \quad \badgetype{System Design} \quad \badgefreq{\freqmed}

\stronganswer
This is a classic case where automated metrics diverge from user experience.  I would
investigate along several axes.

\textbf{1. Input distribution shift:}  Compare the input distribution from the past week
to the baseline period.  Check average prompt length, topic distribution (via clustering
input embeddings), language mix, and any new input patterns.  Users may be asking
different types of questions that the model handles poorly, even though the model itself
has not changed.

\textbf{2. Metric blind spots:}  Automated metrics may not capture what users care about.
If we are monitoring perplexity or BLEU, these correlate poorly with user-perceived quality
for open-ended generation.  Check: are we measuring factuality?  Helpfulness?  Tone?
Add LLM-as-judge evaluation on a sample of recent traffic.

\textbf{3. System-level issues:}  Check for silent failures: increased timeout rates
(responses truncated), load balancer changes routing to a different model version,
cache corruption returning stale responses, or infrastructure changes affecting context
length limits.

\textbf{4. Upstream changes:}  Check for changes in the retrieval system (if RAG),
system prompt, safety filters, or post-processing logic.  Any of these can degrade
output quality without the LLM itself changing.

\textbf{5. User expectations:}  Consider whether user expectations have shifted (e.g.,
competing products improved, creating a relative quality perception change).

\textbf{Diagnosis process:}  Sample 50--100 recent user complaints.  Manually evaluate
them alongside model outputs from both periods.  Categorize failure modes (factual errors,
formatting issues, refusals, irrelevant responses).  This often immediately reveals the
root cause.

\redflags
\begin{itemize}
    \item Assuming the model must have changed if quality degraded.
    \item Not considering input distribution shift.
    \item Relying solely on automated metrics without manual inspection.
    \item Not considering system-level or upstream changes.
\end{itemize}

\interviewtest{Production debugging mindset---can the candidate systematically reason
about failure modes in a complex system beyond just the model?}

\goodgreat
\begin{itemize}
    \item \badgeLfive{L5}: Checks input distribution, model version, and automated
          metric coverage.  Samples complaints manually.
    \item \badgeLsix{L6}: Considers the full system (retrieval, prompts, filters, infra),
          proposes specific monitoring additions, and designs a structured diagnosis
          workflow.
    \item \badgeLseven{L7}: Designs a comprehensive observability stack with
          embedding-space drift detection, LLM-as-judge pipelines, user feedback
          integration, and automatic alerting on quality segments (per language,
          per topic, per user cohort).
\end{itemize}

\followups{How would you set up monitoring to catch this proactively?  How do you
measure ``quality'' for open-ended generation?  What is the role of human evaluation
in production systems?}
\end{interviewq}

\vspace{0.5em}

% --- Question 10: Model Routing ---
\begin{interviewq}
\textbf{Q: Describe a model routing strategy that reduces serving cost by 5$\times$ while
maintaining quality.}

\badgeLsix{L6} \quad \badgetype{System Design} \quad \badgefreq{\freqlow}

\stronganswer
The key insight is that most production queries are ``easy'' and do not require a
frontier model.  A model routing system directs each query to the cheapest model capable
of handling it well.

\textbf{Architecture:}
\begin{enumerate}
    \item A lightweight \textbf{router} (small classifier or embedding model) evaluates
          each incoming query.
    \item \textbf{Easy queries} (simple factual QA, routine classification, formatting
          requests) route to a small model (e.g., 7B INT4)---$\sim$$50\times$ cheaper
          than GPT-4.
    \item \textbf{Hard queries} (complex reasoning, multi-step tasks, ambiguous requests)
          route to a large model (e.g., 70B FP16 or API call to frontier model).
    \item If 80\% of traffic is routable to the small model (typical for many workloads),
          effective cost is: $0.8 \times c_{\text{small}} + 0.2 \times c_{\text{large}}$.
          With $c_{\text{small}} = c_{\text{large}} / 50$, this is $0.016c + 0.2c =
          0.216c$---a $\sim$$4.6\times$ cost reduction.
\end{enumerate}

\textbf{Router design:}
\begin{itemize}
    \item Train a difficulty classifier on a dataset of (query, difficulty) labels, where
          difficulty is determined by comparing small-model and large-model outputs.
    \item Features: query length, embedding, presence of reasoning keywords, topic.
    \item Alternatively: always try the small model first; if the response confidence is
          low or fails a quality check, fall back to the large model.
\end{itemize}

\textbf{Quality guarantee:}  Monitor per-route quality metrics.  If the small model's
quality degrades on any query segment, the router thresholds can be adjusted dynamically
to route more traffic to the large model.

\redflags
\begin{itemize}
    \item Proposing only a single model without considering routing.
    \item Not quantifying the cost savings.
    \item Not addressing how to maintain quality for routed queries.
\end{itemize}

\interviewtest{System-level cost optimization thinking and practical deployment
architecture.}

\goodgreat
\begin{itemize}
    \item \badgeLfive{L5}: Describes the basic routing concept and gives a cost
          calculation.
    \item \badgeLsix{L6}: Designs the router architecture, discusses fallback strategies,
          and proposes quality monitoring per route.
    \item \badgeLseven{L7}: Discusses cascading routers (multiple tiers), learned routing
          policies, the interaction between routing and batching efficiency, and A/B
          testing the router itself.
\end{itemize}

\followups{How would you train the difficulty classifier?  What if the small model fails
silently (confidently wrong)?  How does routing interact with latency SLAs?}
\end{interviewq}


% ============================================================================
% CHAPTER SUMMARY
% ============================================================================
\section*{Chapter Summary}

\begin{enumerate}
    \item \textbf{Tokenization} determines vocabulary size, sequence length, and
          multilingual performance.  BPE (bottom-up merging by frequency), WordPiece
          (merging by likelihood), and Unigram (top-down pruning by likelihood loss)
          are the three core algorithms.  Byte-level BPE eliminates UNK tokens entirely.

    \item \textbf{The serving primitives} you depend on---KV cache (memory
          $2 \times L \times H_{\text{kv}} \times S \times d_h \times b$ per request),
          GQA/MLA, continuous batching, PagedAttention, prefix caching, speculative
          decoding, TP/PP---each buy a specific thing and each impose a specific
          application constraint.  Their mechanics are canonical in [DL~19];
          what belongs to you is prompt length and prompt stability, which set
          prefill cost and prefix-cache hit rate.

    \item \textbf{Compression is a three-way choice}: quantize (no retraining,
          $2$--$4\times$, the default first move), distill (needs a training run, but
          the only lever that changes architecture), or prune (needs hardware support
          to become speed).  Methods are canonical in [DL~14]; the application-side
          rule is to evaluate the compressed model on \emph{your} task, since
          benchmark-average quality loss hides task-specific collapse.

    \item \textbf{Cost is an identity, not a mystery}: prompt tokens $\times$ input
          price plus output tokens $\times$ output price, over every model call the
          request makes.  Routing, semantic caching, prefix-cache-friendly prompt
          ordering, prompt compression, output-length control, and batch tiers are the
          levers---and the tokenizer tax multiplies the whole bill for non-English users.

    \item \textbf{Multilingual deployment} requires awareness of the tokenizer tax,
          cross-lingual transfer limitations, and per-language quality monitoring.

    \item \textbf{Operating an LLM feature} means two latency clocks (TTFT and
          inter-token latency) with separate fixes, quality monitoring assembled from
          implicit signals, sampled automated judging, and a golden set, every metric
          sliced by language---and versioning the whole deployed tuple (model, tokenizer,
          prompt, decoding parameters, tool schemas, index version), not just the model.
          The general lifecycle---drift, retraining, CI/CD, canaries---is [DL~22].
\end{enumerate}

\begin{keyinsight}[Chapter 12 Summary: What Interviewers Are Really Testing]
When interviewers ask production NLP questions, they are testing three things:
(1)~\emph{Do you think in numbers?}---Can you estimate KV-cache memory, prefill latency,
and throughput for a specific model?  Candidates who reason quantitatively signal that
they have actually deployed models, not just trained them.
(2)~\emph{Do you know the full stack?}---Tokenization, serving, cost, and monitoring are
all part of production NLP.  Depth in any one of these is valuable, but breadth across all
of them signals a staff-level engineer.
(3)~\emph{Can you make tradeoffs?}---Every optimization trades something (quality, cost,
complexity, generality).  The best candidates do not just list techniques---they reason
about when each technique is appropriate and when it is not.
\end{keyinsight}
